\section{Parameters, computation, and the final promise}\label{sec:completion}

We now put the constructions in their algorithmic order. This matters because the tail formulas are proved for tuned wells, whereas the positions must be computed before the tuning is performed. The surrogate fields in \cref{sec:hopping} remove this apparent dependence.

\subsection{A uniform choice of scales}

Let $n\ge2$ bound the source input size, with the fixed polynomial magnitude and inverse-gap bounds from \cref{thm:source}. Write $E_{\rm in}=E(H_{\rm in})$. Put
\[
\delta_s=b_s-a_s,\qquad \gamma=\min\{1,\delta_s\},
\qquad \alpha=\eta^2/U.
\]
Use the scales and rational hopping approximation of \cref{prop:hubbard}. They give
\begin{equation}\label{eq:finite-final-error}
\left|E_{\rm Hub}(t)-\alpha(E_{\rm in}-C_s)\right|
\le \frac{3\alpha\gamma}{64},
\qquad
C_s=3k\Delta+\frac32\sum_e v_e^2+
                      \sum_{\{i,j\}\in\mathcal C}K_{ij}.
\end{equation}
All parameters chosen so far have polynomial magnitude or inverse-polynomial lower bounds as appropriate. Their binary descriptions have polynomial length. The scale $\alpha$ is rational. The scalar $C_s$ is a sum of polynomially many efficiently computable square roots and rational terms; it does not depend on $E_{\rm in}$.

Set the requested continuum error to
\begin{equation}\label{eq:continuum-budget}
\epsilon_* = \frac{\alpha\gamma}{64}.
\end{equation}
We verify that a single deterministic polynomial choice of $D$ meets this budget. The calibrated displacement bound has the form
\begin{equation}\label{eq:uniform-displacement-size}
S\le C n^b(1+\log D)
\end{equation}
for a fixed $b$, since $m$, the number of contacts, and $|\log t_{ij}|$ are polynomially bounded in $n$. Increase $b$ so that $m\le Cn^b$ as well. The hopping error before polynomial norm conversions is bounded by a constant times
\begin{equation}\label{eq:calibration-budget}
\begin{split}
\mathfrak e(D,\tau)={}&
\frac{m\log D+m^2}{D}+\frac{mS}{D}
+\frac{S^2}{D}+S D^{-1/4}\\
&+\tau+e^{-c\sqrt D}.
\end{split}
\end{equation}
The other errors are $\poly(m)/D$ and the exponentially small terms listed below. Choose fixed exponents $h,r$ so that all polynomial norm-conversion and coefficient-counting losses are bounded by $Cn^h$, and $\epsilon_*\ge c n^{-r}$. In particular, take $m^2+m^5\le Cn^h$; the two algebraic error terms in \cref{eq:continuum-simulation-final} are then bounded by $Cn^h(\mathfrak e(D,\tau)+D^{-1})$. Such exponents exist by the explicit finite construction. Increasing them only weakens the conditions to be met.

Choose an integer $K$ satisfying
\begin{equation}\label{eq:choice-K}
K>4(b+h+r+2),
\end{equation}
and put
\begin{equation}\label{eq:choice-D}
D=A(n+2)^K,
\qquad
\tau=c_*(n+2)^{-h-r-2},
\end{equation}
where $A$ is a sufficiently large fixed integer and $c_*>0$ a sufficiently small fixed rational. The factors $1+\log D$ in \cref{eq:uniform-displacement-size} do not make this choice circular: $K$ is chosen first, and then $A$ is fixed. Each $D$-dependent expression tends to zero relative to its error budget as $A$ increases, even with its powers of $\log A$. The remaining growth in $n$ is controlled by \cref{eq:choice-K}. In particular,
\[
n^h S D^{-1/4}=o(n^{-r}),\qquad
n^h S^2/D=o(n^{-r}),\qquad
m+S=o(D^{1/4}).
\]
Here the little-oh statements describe the uniform large-scale estimates used to choose the fixed constants; increasing $A$ covers every finite small-$n$ case. The same choices bound all the terms in \cref{eq:calibration-budget} and the Coulomb-compression error by their assigned fractions of \cref{eq:continuum-budget}. No search for a successful $D$ is part of the algorithm.

It is important to keep the energy units explicit. After dilation and subtraction of the isolated-well ground energy,
\begin{equation}\label{eq:units-final}
\mathcal H=H/Z_*^2+m/2,
\qquad Z_*\mathcal H=H/Z_*+mZ_*/2.
\end{equation}
\Needspace{10\baselineskip}
The exponential margins that are needed after multiplication by $Z_*\asymp e^D$ are:
\begin{center}
\begin{tabular}{@{}lll@{}}
\toprule
Source of error & In $\mathcal H$ units & In $Z_*\mathcal H$ units\\
\midrule
Tuned well energies & $m e^{-1.01D}$ & $O(m e^{-.01D})$\\
Discarded one-body entries & $O(m^2 e^{-1.1D})$ & $O(m^2 e^{-.1D})$\\
Orthonormalization & $\poly(m)e^{-1.6D}$ & $\poly(m)e^{-.6D}$\\
Continuum complement & $e^{-1.2D}$ & $O(e^{-.2D})$\\
\bottomrule
\end{tabular}
\end{center}
The polynomial multiplicities are retained in \cref{eq:continuum-simulation-final}. The positive exponential margins allow all these terms to fit their assigned error budgets by increasing the fixed $A$. The one-body residual $e^{-.75D}$ is used quadratically against a fixed complementary gap; multiplying that residual linearly by $Z_*$ would not be a valid estimate. \Cref{prop:continuum-simulation} supplies the required quadratic bound.

With these choices, that proposition gives
\begin{equation}\label{eq:full-final-error}
\left|E_0/Z_*+mZ_*/2-\alpha(E_{\rm in}-C_s)\right|
\le \frac{\alpha\gamma}{16}
\le \frac{\alpha\delta_s}{16}.
\end{equation}

\subsection{The deterministic construction and its bit cost}

First check the source encoding and the fixed polynomial bounds in \cref{thm:source}, including the inverse-gap bound. If a check fails, output the fixed instance with one electron, one unit-charge nucleus at the origin, and thresholds $-1/2$ and $1/2$. Otherwise carry out the following steps. The guard is polynomial-time and makes the reduction a total polynomial-time map; it never rejects a promised source instance.

\begin{enumerate}[label=\arabic*.,leftmargin=*]
\item \emph{Finite models.} Read the signed rational Heisenberg instance. Form the ancillary graph, the algebraic reference positions $p_i$, the complete contact set $\mathcal C$, the explicit scales $\Delta,\eta$, and the positive rational hoppings $t_{ij}$ of \cref{prop:hubbard}. The finite contact classification avoids deciding arbitrary geometric equalities numerically. All square roots here are approximated with certified rational intervals to the stated inverse-polynomial errors.

\item \emph{Charge scale.} Choose $D$ as in \cref{eq:choice-D} and compute a positive integer $Z_*$ with
\[
e^D/2\le Z_*\le2e^D.
\]
This does not require an exact comparison with a transcendental number. For example, compute the rational sum $s_D=\sum_{j=0}^{4D}D^j/j!$ and take $Z_*=\lfloor s_D\rfloor$. The ratio of successive omitted terms is at most $1/4$, and $j!\ge(j/e)^j$, so $0<e^D-s_D<1$ for sufficiently large $D$. Thus this choice has the required bounds. The sum has polynomially many terms and polynomial bit cost. Round $Z_*(1-500m/D)$ and $Z_*1000m/D$ to the positive integer primary and secondary charges. They are fixed before any energy tuning.

\item \emph{Positions.} At the reference arrangement $X_i^0=Dp_i$, compute each positive surrogate hopping integral to the required relative accuracy, then compute
\[
d_{ij}=\log\bigl(Z_*T_{ij}(X^0)/t_{ij}\bigr)
\]
to the prescribed inverse-polynomial error. Apply the explicit right inverse from \cref{lem:displacement} and round the resulting scaled centers to rational $X_i$. The final displacements are defined by $s_i=X_i-Dp_i$; thus the exact contact-vector decomposition is retained after rounding. \Cref{prop:calibration} proves that these operations take polynomial time and satisfy the required projected-displacement tolerances. The calculation uses $W_i^*$ and is independent of the unknown tuned parameters.

\item \emph{Well tuning.} With these final rational centers and integer charges fixed, apply the certified one-well oracle and the contraction iteration of \cref{prop:tuning}. This produces rational $\rho_i\in[1,4]$ such that $|\lambda_i+1/2|\le e^{-1.01D}$. No subsequent change to a charge or center is made.

\item \emph{Physical instance.} Output the $2m$ nuclei at
\[
R_{i,0}=X_i/Z_*,\qquad
R_{i,1}=(X_i+\rho_i e_z)/Z_*
\]
with their positive integer charges, and output $N=m$ in unary. Compute the rational thresholds specified in the next subsection.
\end{enumerate}

Each step is classical and deterministic. The one-well Ritz spaces have dimension polynomial in $D$ and in the rational coordinate bit lengths. Their matrix entries and conditioning bounds require only polynomially many bits, and the number of eigenvalue evaluations in the tuning procedure is polynomial. The integral evaluations and elementary functions in calibration likewise require only polynomial time and precision. In particular, no many-electron ground energy is queried.

The output has polynomial length. Each nuclear charge has $O(D+\log n)$ bits. Rational scaled centers and secondary offsets have polynomial bit length, and division by the $O(D)$-bit integer $Z_*$ preserves that property. Distinct sites are separated by $D-O(S)$ in scaled coordinates; each secondary nucleus is between one and four units from its own primary. Thus all output nuclei are distinct for the chosen $D$. Reducing rational coordinates and thresholds to lowest terms is a polynomial-time integer computation. Finally, $m$ is polynomial in $n$, so unary encoding of $N$ has polynomial length as well. Auxiliary caps, orbitals, Ritz bases, and projections are used in the proof and computation only; none is part of the output instance.

\subsection{Rational thresholds and preservation of both promises}

Compute a rational $\widetilde C_s$ with
\begin{equation}\label{eq:scalar-precision}
|\widetilde C_s-C_s|\le\delta_s/64.
\end{equation}
Polynomially many square-root approximations suffice, since $C_s$ has polynomial magnitude and $\delta_s$ has an inverse-polynomial lower bound. Define the rational affine map
\begin{equation}\label{eq:affine-map}
\widetilde F(z)=-mZ_*^2/2+Z_*\alpha(z-\widetilde C_s),
\end{equation}
and set
\begin{equation}\label{eq:output-thresholds}
a=\widetilde F(a_s+\delta_s/4),
\qquad b=\widetilde F(b_s-\delta_s/4).
\end{equation}
All the quantities in these formulas are rational and computable with polynomial bit cost. The possibly enormous common offset $-mZ_*^2/2$ is an exact rational with polynomial bit length.

By \cref{eq:full-final-error,eq:scalar-precision},
\begin{equation}\label{eq:rounded-affine-error}
|E_0-\widetilde F(E_{\rm in})|
\le \frac{5}{64}Z_*\alpha\delta_s
<\frac18Z_*\alpha\delta_s.
\end{equation}
Since $\widetilde F$ has positive slope, a YES source instance satisfies
\[
E_0\le\widetilde F(a_s)+\tfrac18Z_*\alpha\delta_s
\le a-\tfrac18Z_*\alpha\delta_s<a,
\]
and a NO source instance satisfies
\[
E_0\ge\widetilde F(b_s)-\tfrac18Z_*\alpha\delta_s
\ge b+\tfrac18Z_*\alpha\delta_s>b.
\]
These inequalities concern the exact continuum spectral infimum by \cref{prop:continuum-simulation}, rather than only a projected eigenvalue.

The threshold separation is exactly
\begin{equation}\label{eq:output-gap}
b-a=\tfrac12 Z_*\alpha\delta_s.
\end{equation}
The reciprocal of $\alpha\delta_s$ is polynomially bounded in $n$, whereas $Z_*\ge e^D/2$. Increasing the same fixed scale constant if necessary makes \cref{eq:output-gap} at least one for every input. Since the output length $L\ge1$, this implies $b-a\ge L^{-1}$ without any circular estimate involving $L$.

We have constructed a legal nuclear instance in deterministic polynomial time, preserving YES and NO and permitting all spin sectors and all continuum states. This proves \cref{thm:main}.
